\section{Entailment}

We here discuss entailment in propositinal logics based on contractions of formulas.

%% From report
Entailment is the central consequence relation among logical formulas.
Let us define this relation first based on the models of a knowledge base and a test formula.

\begin{definition}[Entailment of propositional formulas]
    \label{def:logicalEntailment}
    Given two propositional formulas $\kb$ and $\exformula$ we say that $\kb$ entails $\exformula$, denoted by $\kb\models\exformula$, if any model of $\kb$ is also a model of $\exformula$, that is
    \begin{align*}
        \uniquantwrtof{\shortatomindicesin}{\imppremhead{\kbat{\indexedshortcatvariables}=1}{\formulaat{\indexedshortcatvariables}=1}} \, .
    \end{align*}
    If $\kb\models\lnot\exformula$ holds, we say that $\kb$ contradicts $\exformula$.
\end{definition}

% Connection with tensor formalism
To use the tensor network formalism for the decision of entailment, we will in the following develop three equivalent criteria for entailment. %, based on contractions.

\subsect{Deciding Entailment by Contractions}

First of all, we can decide entailment based on vanishing contractions with the negated test formula.

\begin{theorem}[Contraction Criterion of Entailment]
    \label{the:contCriterionLogEntailment}
    We have $\kb\models\exformula$ if and only if
    \begin{align*}
        \contraction{\kb,\lnot\exformula} = 0 \, .
    \end{align*}
\end{theorem}
\begin{proof}
    \proofleftsymbol:
    If for a $\shortatomindicesin$ we have $\kbat{\indexedshortcatvariables}=1$ but not $\big(\exformulaat{\indexedshortcatvariables}=1\big)$, we would have $\big(\lnot\exformulaat{\indexedshortcatvariables}=1\big)$ and
    \begin{align*}
        \contraction{\kb,\lnot\exformula} =
        \sum_{\shortatomindicesin} \kbat{\indexedshortcatvariables} \cdot \exformulaat{\indexedshortcatvariables} > 1 \, .
    \end{align*}
    Thus, whenever the contraction vanishes, we have
    \begin{align*}
        \uniquantwrtof{\shortatomindicesin}{\imppremhead{\kbat{\indexedshortcatvariables}=1}{\formulaat{\indexedshortcatvariables}=1}} \, .
    \end{align*}
    %for all $\shortatomindicesin$ that $\big(\kbat{\indexedshortcatvariables}=1\big) \rightarrow \big(\formulaat{\indexedshortcatvariables}=1\big)$.

    \proofrightsymbol:
    Conversely, if the contraction $\contraction{\kb,\lnot\exformula}$ does not vanish, we would find $\shortatomindicesin$ with $\kbat{\indexedshortcatvariables}=1$ and $\lnot\formulaat{\indexedshortcatvariables}=1$, therefore $\formulaat{\indexedshortcatvariables}=0$.
    It follows that $\kb\models\exformula$ does not hold.
\end{proof}


% Can use basis encoding
The contraction criterion can be extended to the decision of contradiction as well, since $\kb\models\lnot\exformula$ is equivalent to $\contraction{\kb,\exformula}=0$.
Therefore, entailment and contradiction can be decided simultaneously by a single contraction, as we state next.
%To decide whether a formula is entailed, or its negation is entailed (in which case one says that the formula is contradicted) by a single contraction, one can perform the contraction
%\begin{align*}
%	\hypercore = \contractionof{\kbat{\shortcatvariables},\formulaat{\shortcatvariables,\exformulavar}}{\exformulavar}
%\end{align*}

\begin{theorem}
    \label{the:entailmentContradictionContraction}
    Given propositional formulas $\kb$ and $\exformula$ we build
    \begin{align*}
        \hypercoreat{\exformulavar}
        = \contractionof{\kbat{\shortcatvariables},\formulaccwith}{\exformulavar} \, .
    \end{align*}
    Then $\kb\models\exformula$ is equivalent to $\hypercoreat{\exformulavar=0}=0$, and $\kb\models\lnot\exformula$ is equivalent to $\hypercoreat{\exformulavar=1}=0$ \, .
\end{theorem}
\begin{proof}
    This follows from \theref{the:contCriterionLogEntailment} using that
    \begin{align*}
        \contraction{\kb,\lnot\exformula} = \hypercoreat{\exformulavar=0}
    \end{align*}
    and
    \begin{align*}
        \contraction{\kb,\exformula} = \hypercoreat{\exformulavar=1} \, . & \qedhere
    \end{align*}
\end{proof}





\subsect{Deciding Entailment by Partial Ordering}

% Subset relation
Logical entailment can be understood by subset relations of the models of the respective formulas.
This perspective can be applied with subset encodings in \charef{cha:basisCalculus}.
The subset relation corresponds with partial ordering of its encoded tensors, as will be shown in \theref{the:subsetRelationSubsetEncoding}.
For two propositional formulas, we denote to this end $\exformula\prec\secexformula$ (see \defref{def:partialOrder}), if and only if for all $\shortatomindicesin$
\begin{align*}
    \exformulaat{\indexedshortcatvariables} \leq \secexformulaat{\shortcatindices}  \, .
\end{align*}

\begin{theorem}[Partial Ordering Criterion of Entailment]
    \label{the:orderingEntailmentCriterion}
    We have $\kb\models\exformula$ if and only if $\kbat{\shortcatvariables}\prec\exformulaat{\shortcatvariables}$.
\end{theorem}
\begin{proof}
    Since both $\kb$ and $\exformula$ are boolean tensors, we have for any $\shortatomindicesin$ that
    \begin{align*}
        \kbat{\indexedshortcatvariables},\formulaat{\indexedshortcatvariables}\in\ozset \, .
    \end{align*}
    Thus,
    \begin{align*}
        \uniquantwrtof{\shortatomindicesin}{\kbat{\indexedshortcatvariables}\leq\formulaat{\indexedshortcatvariables}}
    \end{align*}
    is equivalent to
    \begin{align*}
        \uniquantwrtof{\shortatomindicesin}{\imppremhead{\kbat{\indexedshortcatvariables}=1}{\formulaat{\indexedshortcatvariables}=1}} \, .
    \end{align*}
    This states that $\kbat{\shortcatvariables}\prec\exformulaat{\shortcatvariables}$ is equivalent to $\kb\models\exformula$.
\end{proof}



\subsect{Redundancy of Entailed Formulas}

Another interpretation of entailment is by redundancy of a formula in a Knowledge Base.
This is especially interesting for the sparse representation of Knowledge Bases.
%Towards getting insights on this we first show that entailed formulas can be dropped from the Knowledge Base.

\begin{theorem}[Redundancy Criterion of Entailment]
    \label{the:ReduncancyOfEntailed}
    If and only if $\kb\models\exformula$ we have
    \begin{align*}
        \kbat{\shortcatvariables}= \contractionof{\kb,\exformula}{\shortcatvariables}  \, .
    \end{align*}
\end{theorem}
\begin{proof}
    For any formula $\exformula$ we have
    \begin{align*}
        \onesat{\shortcatvariables} = \exformulaat{\shortcatvariables} + \lnot\exformulaat{\shortcatvariables}
    \end{align*}
    and thus
    \begin{align*}
        \kbat{\shortcatvariables}
        & = \contractionof{\kbat{\shortcatvariables},\onesat{\shortcatvariables}}{\shortcatvariables} \\
        & = \contractionof{\kbat{\shortcatvariables},\exformulaat{\shortcatvariables}}{\shortcatvariables} +  \contractionof{\kbat{\shortcatvariables},\lnot\exformulaat{\shortcatvariables}}{\shortcatvariables} \, .
    \end{align*}
    Now, by \theref{the:contCriterionLogEntailment} we have $\kb\models\exformula$, if and only if $\contractionof{\kbat{\shortcatvariables},\lnot\exformulaat{\shortcatvariables}}{\shortcatvariables}=0$, which is thus equal to
    \begin{align*}
        \kbat{\shortcatvariables}
        = \contractionof{\kbat{\shortcatvariables},\exformulaat{\shortcatvariables}}{\shortcatvariables} \, . & \qedhere
    \end{align*}
\end{proof}

%This provides us with another interpretation of the entailment relation, in terms of redundant formulas.
%\begin{remark}[Sparsest Description of a Knowledge Base]
%	Given a set of worlds indexed by $\hypercore$, find the sparsest set of formulas $\kb$ such that
%		\[ \hypercore = {\kb} \]
%	would be benefitial for small computational complexity.
%	Since the formula tensors are invariant under entailment, we can drop entailed formulas.
%\end{remark}

\subsect{Contraction Knowledge Base}

We exploit the contraction and redundancy criteria of entailment to sketch an implementation of a propositional Knowledge Base in \algoref{alg:contractionKB}.
Here the function $\mathrm{ASK}(\exformula)$ returns, whether a formula $\exformula$ is entailed or contradicted by a Knowledge Base.
If the formula is neither entailed or contradicted, we say it is contingent.
If it is both, we have $\kbat{\shortcatvariables}=0$ and thus an inconsistent Knowledge Base.
Exploiting \theref{the:entailmentContradictionContraction} we decide these situations based on a single contraction.

The function $\mathrm{TELL}(\exformula)$ incorporates an additional formula $\exformula$ into a Knowledge Base $\kb$.
Here we exploit \theref{the:ReduncancyOfEntailed} and do not add a formula, which is entailed in order to maintain a sparse representation. %(in which case it returns Incons).
The function further refuses to add a formula, which would make the Knowledge Base inconsistent (returns Refused) and only changes the Knowledge Base in case of a contingent formula (returns Added).


\begin{figure}
    \begin{center}
        \input{partI/0-propositional-logics/tikz/entailment/ask-decision-table.tex}
    \end{center}
    \caption{Table of possible logical relations between a knowledge base $\kb$ and $\formula$, based on whether the knowledge base entails the formula ($\kb\models\formula$) and its negation ($\kb\models\lnot\formula$).}\label{fig:askDecisionTable}
\end{figure}


% Works also for Markov Networks!
\begin{algorithm}[hbt!]
    \caption{Contraction Knowledge Base with operations ASK and TELL}\label{alg:contractionKB}
    \vspace{0.3cm} \textbf{ASK}($\kb$, $\exformula$)
    \begin{algorithmic}
    \iosepline
    \Require Knowledge base $\kb$, query formula $\exformula$
    \Ensure Decision which relation between $\kb$ and $\exformula$ holds (see \figref{fig:askDecisionTable})
    \iosepline
        \State{$\hypercoreat{\formulavar} \algdefsymbol \contractionof{\{\secexformulaat{\shortcatvariables} \wcols \secexformula\in\kb\},\bencodingofat{\exformula}{\formulavar,\shortcatvariables}}{\formulavar}$}
        \If{$\hypercoreat{\formulavar=0}=0$ and $\hypercoreat{\formulavar=1}=0$ }
            \State \Return \stringof{Inconsistent}
        \ElsIf{$\hypercoreat{\formulavar=0}=0$}
            \State \Return \stringof{Entailed}
        \ElsIf{$\hypercoreat{\formulavar=1}=0$}
            \State \Return \stringof{Contradicted}
        \Else
            \State \Return \stringof{Contingent}
        \EndIf
    \vspace{0.1cm}
    \hrule
    \end{algorithmic}
    \vspace{0.3cm} \textbf{TELL}($\kb$, $\exformula$)
    \begin{algorithmic}
        \iosepline
    \Require Knowledge base $\kb$, query formula $\exformula$
    \Ensure Decision whether a formula is added to the knowledge base (\stringof{Added}), or an exception in (\stringof{Inconsistent}, \stringof{Redundant} or \stringof{Refused}) is raised.
    \iosepline
        \State{answer $\algdefsymbol$ ASK($\exformula$)}
        \If{answer is \stringof{Inconsistent}:}
            \State \Return \stringof{Inconsistent}
        \ElsIf{answer is \stringof{Entailed}:}
            \State \Return \stringof{Redundant}
        \ElsIf{answer is \stringof{Contradicted}:}
            \State \Return \stringof{Refused}
        \ElsIf{answer is \stringof{Contingent}:}
            \State $\kb \algdefsymbol \kb\cup\{\exformula\}$
            \State \Return \stringof{Added}
        \EndIf
    \end{algorithmic}

\end{algorithm}